\section{Lời giải thích do chính mô hình viết}
\label{sec:ch17-loi-giai-thich-tu-viet}

\index{lời giải thích tự sinh!khác gì lời giải thích hậu nghiệm|(}%
Phần mở đầu chương vừa nói vế trái của phép so sánh đổi chủ. Mục này nói đổi chủ
thì kéo theo những gì, vì bốn hệ quả của nó là bốn thứ mà sáu mục sau đều dựa
vào.

\index{lời giải thích tự sinh!định nghĩa dùng trong chương}%
Trước hết phải đặt tên cho thứ đang bàn. Một
\tn{lời giải thích tự sinh}{self-explanation}, trong chương này, là đoạn văn bản
mô tả một quá trình mà chính mô hình in ra cùng với đáp án, chứ không phải do
một thủ tục bên ngoài dựng lên sau khi đáp án đã có. \tn{Chuỗi suy luận}{chain-of-thought}
mà mục~\ref{sec:ch02-chuoi-suy-luan} mô tả là dạng thường gặp nhất của nó, và là
dạng duy nhất chương này đọc. Chữ \enquote{tự sinh} nói về nguồn gốc của đoạn
văn bản chứ không nói nó đúng: một lời giải thích tự sinh vẫn có thể mô tả một
quá trình chưa từng xảy ra, và mục~\ref{sec:ch17-bai-neo-tro-ve} tính ra một ví
dụ như vậy.

Hình~\ref{fig:ch17-hai-nguon-cua-ve-trai} đặt hai nguồn ấy cạnh nhau. Bên trái là
hình dạng của các phương pháp mà Phần~II dựng: LIME, SHAP, Integrated Gradients
và Grad-CAM đều là một thủ tục chạy riêng, nhận mô hình và đầu vào rồi trả về một
lời giải thích. Bên phải là hình dạng của chương này. Hai chương của sách không
nằm gọn ở cột nào: chương~\ref{ch:chi-so-khong-do-do-trung-thuc} đọc một bài mà
vế trái đã là chuỗi suy luận, tức đã ở cột phải, và
chương~\ref{ch:nut-that-khai-niem} đọc những mô hình mà vế trái là một tầng của
chính mô hình chứ không phải đầu ra của thủ tục nào. Chỗ đứng lệch của chương
đầu là chuyện mà mục~\ref{sec:ch17-chieu-cua-phep-chuyen} quay lại.

\begin{figure}[htbp]
  \centering
  \begin{tikzpicture}[
  every node/.style={font=\footnotesize, align=center},
  known/.style={draw=black!85, line width=1pt, fill=black!12,
    minimum width=32mm, minimum height=9mm},
  unknown/.style={draw=black!55, dotted, thick, fill=white,
    minimum width=38mm, minimum height=9mm},
  route/.style={draw=black!65, rounded corners=2pt, fill=black!5,
    minimum width=40mm, minimum height=20mm, text width=38mm,
    font=\scriptsize},
  flow/.style={-{Stealth[length=2mm]}, thick, black!70},
  thinflow/.style={-{Stealth[length=1.6mm]}, black!55},
  lbl/.style={font=\scriptsize, align=center, black!55}
]
  \node[known]   (gt)  at (1.7,4.3) {vế trái:\\lời giải thích};
  \node[unknown] (dep) at (8.3,4.3) {vế phải: cái mà hành vi\\phụ thuộc vào};

  \draw[flow] (gt.east) -- (dep.west);
  \node[lbl] at (5.0,4.95) {định nghĩa~\ref{def:ch01-do-trung-thuc} so hai bên};

  \node[route] (r1) at (0,0.4)
    {\textbf{một phương pháp bên ngoài}\\[1mm]
     nhận mô hình và đầu vào, chạy riêng, rồi tính ra lời giải thích};
  \node[route] (r2) at (5.6,0.4)
    {\textbf{chính mô hình}\\[1mm]
     in chuỗi suy luận và đáp án trong cùng một lượt sinh; không có phương pháp nào chạy riêng};

  \draw[thinflow] (r1.north) -- (0,2.6);
  \draw[thinflow] (r2.north) -- (5.6,2.6);
  \draw[black!55] (0,2.6) -- (5.6,2.6);
  \draw[flow] (1.7,2.6) -- (gt.south);
  \node[lbl] at (2.8,3.15) {hai chỗ vế trái có thể tới từ};

  \node[lbl] at (0,-1.55)   {các phương pháp của Phần~II};
  \node[lbl] at (5.6,-1.55) {chương này};
\end{tikzpicture}

  \caption{Vế phải của định nghĩa~\ref{def:ch01-do-trung-thuc} không đổi giữa
  hai cột, và cả hai cột đều không biết nó. Cái đổi là vế trái: bên trái nó là
  đầu ra của một thủ tục có thể hỏi han được, bên phải nó là một phần đầu ra của
  chính mô hình. Mọi hệ quả trong mục này đến từ chỗ mũi tên bên phải không đi
  qua hộp nào.}
  \label{fig:ch17-hai-nguon-cua-ve-trai}
\end{figure}

\index{lời giải thích tự sinh!bốn hệ quả}%
Hệ quả thứ nhất là không còn tham số nào để hỏi.
Mục~\ref{sec:ch14-dieu-chua-giai-quyet} đặt lên mọi dụng cụ đo tương lai một yêu
cầu, rằng một phương pháp phải nói điểm số của nó tương đối với tập can thiệp
nào, và yêu cầu ấy giả định có một phương pháp để mà hỏi. Ở đây không có. LIME có
bảy lựa chọn tự do mà mục~\ref{sec:ch03-cac-tham-so-tu-do} đếm ra, Integrated
Gradients có baseline mà mục~\ref{sec:ch05-baseline} đọc, SHAP có
\tn{phân phối nền}{background distribution};
một chuỗi suy luận không có cái nào trong số ấy. Điều đó không làm bài toán dễ
đi. Nó bỏ đi luôn cả chỗ để sửa: khi một lời giải thích hậu nghiệm sai, hỏi được
lựa chọn nào đã làm nó sai, còn ở đây không có lựa chọn nào để hỏi.

Hệ quả thứ hai là lời giải thích và cái được giải thích ra đời cùng lúc, từ cùng
một quá trình. Một \tn{bản đồ nhiệt}{heatmap} được tính sau khi mô hình đã dự
đoán xong, nên
việc tính nó không đổi dự đoán. Một chuỗi suy luận thì nằm trong ngữ cảnh mà đáp
án được sinh ra, nên nó vừa có thể mô tả cái đã xảy ra vừa có thể là nguyên nhân
của cái sắp xảy ra. Hai vai ấy không tách được bằng cách đọc chuỗi, và
mục~\ref{sec:ch17-khao-sat-biet-cho-hong} cho thấy một khảo sát năm 2026 gộp
đúng hai vai này lại thành một định nghĩa.

Hệ quả thứ ba là lời giải thích bây giờ là văn bản, nên nó khẳng định được những
điều không có thật. Một vector điểm số không nói được câu nào sai: nó gán cho mỗi
\tn{đặc trưng}{feature} một con số, và con số ấy hoặc phản ánh đúng cái mô hình
dùng hoặc không. Một câu tiếng Anh thì nói được rằng mô hình đã nhớ lại một cuốn sách lịch
sử, và câu ấy sai theo một kiểu mà không họ chỉ số nào của
chương~\ref{ch:do-do-trung-thuc} được thiết kế để bắt.

\index{lời giải thích tự sinh!đối tượng đã đổi một lần}%
Hệ quả thứ tư là hệ quả mà chương này tìm ra và nó nằm ở chỗ khác với ba hệ quả
trên: đối tượng đã đổi một lần rồi, giữa lúc bài mở đường viết ra nó và lúc bài
neo của sách đo nó. Bài neo ghi thẳng điều ấy khi nó viết lại các định nghĩa:
định nghĩa được dùng rộng rãi nhất về độ trung thực được dựng cho những phương
pháp có mục đích sinh ra một lời giải thích mạch lạc và tự chứa về hành vi của
mô hình, ví dụ bản đồ chú ý, các điểm số saliency, và \emph{những chuỗi suy luận
few-shot tuyến tính nguyên bản của Wei và cộng sự}; còn chuỗi suy luận của các mô
hình suy luận hiện nay là những lời nói dài mà mô hình phát ra trong lúc giải một
nhiệm vụ, chứa nhiều bước không nối trực tiếp với dự đoán và có bước không liên
quan gì tới nó~\autocite{p21metrics}.

Câu ấy đáng dừng lại, vì nó nói rằng hai thứ mang cùng một tên không cùng một
loại. Chuỗi của bài mở đường là một lời giải thích theo nghĩa cũ, ngắn, tuyến
tính, viết ra để dẫn tới đáp án, và bài neo xếp nó vào cùng lớp với bản đồ chú ý.
Chuỗi của một mô hình suy luận là dấu vết của việc giải bài, dài, rẽ nhánh, có
đoạn bỏ đi giữa chừng, và bài neo xếp nó ra ngoài lớp ấy.

\index{lời giải thích tự sinh!khác gì lời giải thích hậu nghiệm|)}%
Mục~\ref{sec:ch09-dinh-nghia-lai} đọc chính đoạn văn ấy, và đọc nó để tách chuỗi
suy luận khỏi đối tượng của Phần~II, nên hai thứ đầu trong danh sách là hai thứ
chương ấy cần. Thứ ba mang tải ở đây và không mang tải ở đó, vì nó cho thấy ranh
giới còn chạy cả bên trong chuỗi suy luận. Điều cần giữ từ đây là một câu: khi
hỏi một chuỗi suy luận có trung thực không, phải hỏi thêm chuỗi của loại mô hình
nào, và mục~\ref{sec:ch17-bai-neo-tro-ve} nói chỗ nào trong sách chịu ảnh hưởng
vì câu ấy.
